\documentclass[11pt,a4paper]{article}
\usepackage[a4paper,margin=2.2cm]{geometry}
\usepackage[T1]{fontenc}
\usepackage{lmodern}
\usepackage{microtype}
\usepackage{booktabs}
\usepackage{tabularx}
\usepackage{array}
\usepackage[table]{xcolor}
\usepackage{enumitem}
\usepackage{titlesec}
\usepackage[british]{babel}
\usepackage[colorlinks=true,linkcolor=blue!50!black,urlcolor=blue!50!black]{hyperref}

\setlist{itemsep=0.25ex,topsep=0.4ex,leftmargin=1.4em}
\titlespacing*{\section}{0pt}{2.2ex}{1.0ex}
\titlespacing*{\subsection}{0pt}{1.6ex}{0.6ex}
\setlength{\parindent}{0pt}
\setlength{\parskip}{0.8ex}
\renewcommand{\arraystretch}{1.12}
\newcommand{\iss}[1]{\href{https://github.com/JieGH/RVV_code_gen_via_MLIR_xDSL/issues/#1}{\##1}}
\newcommand{\good}[1]{\textcolor{green!45!black}{\textbf{#1}}}
\newcommand{\slow}[1]{\textcolor{red!70!black}{#1}}

\begin{document}

%-------------------------------------------------------------------- title page
\begin{titlepage}
\centering
\vspace*{2.5cm}
{\LARGE\bfseries Where does the speed go?\par}
\vspace{0.6cm}
{\Large Performance of our generated GEMM design against the reference design on the
SpacemiT X60: findings of issues \#31 to \#43\par}
\vspace{1.2cm}
{\large Engineering report\par}
\vspace{0.4cm}
{\normalsize \texttt{RVV\_code\_gen\_via\_MLIR\_xDSL}, branch \texttt{pure-xdsl}\par}
\vspace{0.2cm}
{\normalsize 28 September 2026\par}
\vfill
\begin{minipage}{0.88\linewidth}\setlength{\parskip}{0.8ex}
\textbf{Summary}\par\smallskip
This report brings together thirteen issues (\#31--\#39, \#41--\#43) that looked into why our design, a matrix
multiplication (GEMM) whose inner kernels are generated with xDSL, is slower on some test layers than expected
and than the reference design. The work ran from 24 to 28 September 2026 on two BananaPi M2 Pro boards.

The main finding is that \textbf{our generated code is not the cause}. At the same tile size, our kernels use the
same number of vector registers as the reference's kernels, spill exactly the same values to memory, and run at
the same speed when one is swapped for the other. The losses come from three other places: tiles that need more
than the 32 vector registers the hardware has; the shape of some layers, which forces more packing work per
operation (layer B4); and a memory ``slow state'' on layer B3, in which the processor's L1 data prefetcher does
only about half its normal work, so that both packing and computing slow down. A first fix for B3, clearing the
pack buffers with the \texttt{cbo.zero} instruction, raises B3 from about 5.6 to about 6.0 GOPs on both boards,
but the reason it helps is not yet understood. The report also shows that measurement noise (between 1 and 5\%
from one launch to the next, and 5--10\% on core 0) was large enough to make several early claims unsafe, and it
records the testing rules adopted as a result.
\end{minipage}
\vspace{1.5cm}
\end{titlepage}

\tableofcontents
\newpage

%-------------------------------------------------------------------- 1
\section{Introduction}

\subsection{Background}
This project generates fast matrix-multiplication code for RISC-V processors with vector units. It follows the
BLIS structure: the large matrices are cut into blocks, each block is copied (``packed'') into a compact buffer,
and a small, highly tuned routine called the \emph{microkernel} then multiplies one small tile of the result at
a time. In our design the microkernels, the packing routines and the loops around them are generated by a
Python/xDSL tool. The \emph{reference design} is an existing, largely hand-written benchmark suite
(\texttt{RVV\_ukernels\_benchmark}) whose microkernels were produced with the Exo tool. It serves as the
baseline for both correctness and speed.

Earlier work brought the two designs close together in method: they use the same blocking rules and, since issue
\#13, both count packing time in their speed figures. Even so, some layers ran more slowly than expected, and
the speed of some layers moved a great deal from one run to the next. The issues covered here were opened to
find out why.

\subsection{Aims}
The work set out to answer four questions:
\begin{enumerate}
  \item Is our code generator (xDSL) producing worse microkernels than the reference?
  \item Does the choice of tile size, in particular tiles that use every vector register, limit speed?
  \item Why does layer B4 run below the square test shapes, even with the same tile?
  \item Why does layer B3 run at two very different speeds, and can it be kept at the faster one?
\end{enumerate}
A fifth aim grew out of the first four: to find out how much of what was measured was real, and how much was
noise.

\subsection{Scope}
Table~\ref{tab:issues} lists the issues. Number \#40 is a merged pull request, not an investigation, and is not
covered. Issues \#37, \#38 and \#44 were written as plans and are on hold; they appear only as recommendations.
The detailed results are in four earlier briefs \cite{b3134,b353639,b42,b43}. This report does not repeat every
number from them. It brings their results together, says what they mean as a whole, and states what is still
unknown.

\begin{table}[h]
\centering\small
\begin{tabularx}{\linewidth}{@{}l X l@{}}
\toprule
\textbf{Issue} & \textbf{Question} & \textbf{State} \\
\midrule
\iss{31} & Does the 32x7 tile (all 32 registers in use) slow B4 down? & closed: no \\
\iss{32} & How do B3 and B4 change as the tile width changes, with each packer? & closed \\
\iss{33} & Does xDSL use registers worse than the reference at 16x16 and 24x10? & closed: no \\
\iss{34} & Why is B4 slower than the square shapes at the same tile? & closed: answered \\
\iss{35} & What does a smaller packed-B block (\texttt{nc}) do for B3? & done; decision made \\
\iss{36} & Why is the vector packer slow on B3? & measured; cause narrowed \\
\iss{37} & 32x6 or 32x7 for B4? & on hold; keep 32x7 \\
\iss{38} & Cheaper A packing and matrix padding for B4 & on hold \\
\iss{39} & Rules for trustworthy benchmark results & adopted \\
\iss{41} & What is B3's slow state? & open; main cause named \\
\iss{42} & Tooling to count L2 and prefetch events per repetition & built; first results \\
\iss{43} & Fix 1: clear pack-buffer lines with \texttt{cbo.zero} & built; measured; off by default \\
\iss{44} & Fix 2: pack B one column at a time & specified; on hold \\
\bottomrule
\end{tabularx}
\caption{Issues covered by this report.}
\label{tab:issues}
\end{table}

\subsection{Terms used}
A few technical terms cannot be avoided. They are used in this report as follows.
\begin{description}[leftmargin=1.2em,itemsep=0.2ex,font=\normalfont\bfseries]
  \item[GOPs] Billions of arithmetic operations per second; higher is faster. Our figures include packing time.
  ``By mean time'' is the reference's formula (all operations over all time); ``by median rep'' uses the middle
  repetition.
  \item[Tile (MR$\times$NR)] The size of the small piece of the result that one microkernel call works on, for
  example 32x7 (32 rows, 7 columns).
  \item[Vector register] A small, very fast store inside the processor. The X60 has 32. A tile that needs more
  than 32 must ``spill'': move values out to memory and back, which is slow.
  \item[Packing] Copying a block of an input matrix into a compact buffer in the order the microkernel reads it.
  The \emph{vector packer} does this with vector instructions, the \emph{scalar packer} one value at a time.
  \item[\texttt{nc}] The width of the packed block of matrix B. Larger values mean fewer, larger blocks.
  \item[Cache, L1, L2] Small, fast memories between the processor and main memory. Each X60 core has a 32\,KiB
  L1 data cache; each group of four cores shares a 512\,KiB L2.
  \item[Prefetcher] A part of the processor that guesses which memory will be needed next and loads it into the
  cache in advance, so that the program does not have to wait for it.
  \item[Launch, rep] A launch is one run of the benchmark program; inside it, each case is repeated many times
  (reps).
  \item[Instrumented build] A build that also reads hardware counters. It is used to see where time goes; its
  speed figures are not compared with the reference.
\end{description}

%-------------------------------------------------------------------- 2
\section{Method}

\subsection{Hardware}
All tests ran on two BananaPi M2 Pro boards, named \texttt{jbpi} and \texttt{jbpi2}. Each has a SpacemiT X60
processor with 8 cores in two groups of four, 256-bit vectors and 32 vector registers.

\subsection{Test cases}
The test cases are matrix multiplications of size M$\times$N$\times$K. The main ones are:
\begin{itemize}
  \item \textbf{S1--S5}: square shapes of increasing size, used as a stable control.
  \item \textbf{B1--B5}: layer shapes from the BERT language model. The two studied most are
  \textbf{B3} (64$\times$384$\times$384, with few rows) and \textbf{B4} (4096$\times$384$\times$1024, with many
  rows but a narrow N). \textbf{B2} (384$\times$384$\times$64) served as a control.
  \item \textbf{Reuse24} (384$\times$384$\times$384) and similar cases, added later to see how the number of rows
  affects B3's behaviour.
\end{itemize}

\subsection{Builds and timing}
Unless stated otherwise, every speed figure in this report comes from the uninstrumented build
(\texttt{INSTRUMENTED=0}) and the reference's own timing method: no warm-up, repetitions run until at least
10 seconds have passed. Figures from the instrumented build are marked and are used only to explain where time
goes. The core each run was fixed to is named where it matters; see Section~\ref{sec:noise} for why.

\subsection{Methods of investigation}
Four methods were used, often together:
\begin{enumerate}
  \item \textbf{Static checks} of the generated code and the compiled machine code: counting live registers,
  spills and instructions, for both designs.
  \item \textbf{Swap tests}: running the reference's microkernel inside our program, so that only the kernel
  differs.
  \item \textbf{Sweeps}: changing one setting (tile width, \texttt{nc}, N, M) while keeping everything else the
  same, ideally inside one launch.
  \item \textbf{Per-rep records}: new tooling (\iss{39}, \iss{42}) saves every repetition's time and, in the
  instrumented build, its hardware counts, so fast and slow repetitions can be compared directly.
\end{enumerate}
Issues \#31 to \#33 were first written up by earlier model versions. They were later checked line by line
against their raw logs \cite{b3134}, and the conclusions below use the checked versions.

%-------------------------------------------------------------------- 3
\section{Results}

\subsection{Our generated code against the reference's code (\#31, \#33)}
The first question was whether xDSL produces worse microkernels. It does not.
\begin{itemize}
  \item \textbf{Registers.} At 16x16 our kernel needs 34 registers and at 24x10 it needs 33, exactly the same as
  the reference's matching kernel style (\texttt{macc}). At 32x7 both need exactly 32.
  \item \textbf{Spills.} After compiling both at \texttt{-O3}, the number of values spilled to memory is the same
  on both sides: 11 or 12 at 16x16 and 4 at 24x10. The instruction counts are within 4\% (ours slightly
  fewer). The extra copy statements visible in our generated source are removed by the compiler.
  \item \textbf{Speed with only the kernel swapped} (Table~\ref{tab:swap}). In every case the difference is
  smaller than the normal change between launches.
\end{itemize}

\begin{table}[h]
\centering\small
\begin{tabular}{@{}llccc@{}}
\toprule
\textbf{Case} & \textbf{Board, core} & \textbf{Our kernel} & \textbf{Reference kernel} & \textbf{Reference} \\
 & & & \textbf{in our program} & \textbf{program} \\
\midrule
16x16, 2048$^3$ & \texttt{jbpi}, 4 & 3.92 & 3.91 & 3.65 (\texttt{macc}) \\
24x10, 2048$^3$ & \texttt{jbpi}, 4 & 6.35 & 6.26 & 6.45 (\texttt{macc}) \\
B4, 32x7 & \texttt{jbpi2}, 4 & 7.72 & 7.74, 7.85 & --- \\
\bottomrule
\end{tabular}
\caption{Kernel swap tests, GOPs by mean time, uninstrumented, one launch per value. The reference's other kernel
style (\texttt{bcast}), which needs 43--50 registers, ran at only 0.90 and 1.45 GOPs.}
\label{tab:swap}
\end{table}

A result from \#31 that our kernel was 7\% faster than the reference's on B4 was measured on core 0 and did not
repeat on core 4. It should not be quoted.

\subsection{Tile size and the 32-register limit (\#31, \#32, \#37)}
Table~\ref{tab:tiles} shows how speed changes with the tile width NR. The rule is simple: \textbf{speed falls
sharply only when a tile needs more than 32 registers}. Below or at the limit, speed changes smoothly, and there
is no sign of a hidden cost for using all 32. The tiles we ship (16x15 for B2 and B3, 32x7 for B4 and the square
shapes) are at or near the best.

\begin{table}[h]
\centering\small
\begin{tabular}{@{}lcccc@{}}
\toprule
\textbf{Case} & \textbf{Shipped tile} & \textbf{GOPs} & \textbf{One column wider} & \textbf{GOPs} \\
\midrule
B2, vector packer & 16x15 (32 regs) & 12.81 & 16x16 (34 regs) & \slow{8.00} \\
B3, vector packer & 16x15 (32 regs) & 5.28 & 16x16 (34 regs) & \slow{3.11} \\
B3, scalar packer & 16x15 (32 regs) & 8.25 & 16x16 (34 regs) & \slow{4.86} \\
B4, vector packer & 32x7 (32 regs) & 7.84 & 32x8 (36 regs) & \slow{4.43} \\
\bottomrule
\end{tabular}
\caption{The register cliff. \#32, \texttt{jbpi}, uninstrumented, GOPs by mean time. These runs were on core 0,
so the absolute values are a little low; the comparison inside each row is still valid.}
\label{tab:tiles}
\end{table}

One open point was whether 32x6 (28 registers) is better than 32x7 for B4. In \#32 it was ahead in all eight
pairs (8.05 against 7.84 GOPs, for example), but all of these ran on core 0. In \#34 on core 4, 32x7 was ahead;
on core 2, 32x6 was ahead by a small amount. The two tiles are therefore level within the noise. \#37 describes a
fair test inside one launch; it is on hold, and 32x7 stays.

\subsection{Why B4 is slower than the square shapes (\#34, \#38)}
B4 runs its compute loop within 2\% of the speed of the square shape S1, so the kernel is not the problem. The
loss is in packing, and it comes from the shape:
\begin{itemize}
  \item B4 has a narrow N (384, like all BERT layers). The blocking rules therefore set \texttt{nc} = 384, and each
  packed block of A is used for only 55 tiles, against about 145 in the square shapes. B4 packs about 2.3 times as
  many values per operation.
  \item In the instrumented build packing takes 14.7\% of B4's time, against about 8\% for S1.
  \item Widening N alone, with M and K unchanged, raises speed from 8.10 GOPs (N=384) to 9.02 GOPs (N=1024), and
  it then stays flat (9.02 at N=2048, 8.97 at N=4096; core 4, one launch). S1 in the same launch ran at 8.73.
  \item A smaller second cost comes from M = 4096 being a power of two: M = 4064 or 4128 runs 2--5\% faster.
\end{itemize}
The reference design packs in the same way and should pay the same cost for this shape. \#38 lists two possible
changes (a cheaper A pack, and padding the matrix row lengths away from powers of two); it is on hold.

\subsection{B3's two speeds (\#32, \#35, \#36, \#41, \#42)}
B3 is the least stable case. The investigation moved in four steps.

\paragraph{Step 1: it looked like a slow packer.} In \#32 the scalar packer beat the vector packer on B3 in every
run (for example 8.25 against 5.28 GOPs at 16x15). On every other shape the vector packer was level or slightly
ahead. Inside one B3 run, speed jumped between about 4.4 and 9.9 GOPs.

\paragraph{Step 2: it is a slow state, not a slow packer (\#36).} The per-rep records showed that every B3 run
starts fast and, after a fraction of a second (vector packer) or up to 8 seconds (scalar packer), drops into a
state about half as fast, where it mostly stays. In that state \emph{both} packing (1.6--3.1 times longer) and
the microkernel (16--66\% longer) slow down, while the clock speed stays the same. The core is waiting on memory.
When B3 is in its fast state, the vector packer is as fast as the scalar one or faster. Both packers read memory
in exactly the same order; the vector one simply sends its requests faster. The clock, large memory pages and B
not fitting in the L2 were all ruled out.

\paragraph{Step 3: a smaller \texttt{nc} helps only in part (\#35).} Table~\ref{tab:nc} shows B3 with different
\texttt{nc} values. With the scalar packer, B3 is steady and fastest at \texttt{nc} = 96--128. With the vector
packer we ship, a smaller \texttt{nc} raises the mean but B3 still has two speeds, until \texttt{nc} is so small
(32) that the packed block fits inside the L1 cache. The blocking model itself cannot produce a small \texttt{nc}
for B3, even when it is given the board's real L2 size.

\begin{table}[h]
\centering\small
\begin{tabular}{@{}r|ccc|ccc@{}}
\toprule
& \multicolumn{3}{c|}{\textbf{Vector packer (shipped)}} & \multicolumn{3}{c}{\textbf{Scalar packer}} \\
\texttt{nc} & by mean time & by median rep & slow reps & by mean time & by median rep & slow reps \\
\midrule
32  & 6.53 & 6.53 & 0\%  & 5.31 & 5.31 & 0\% \\
64  & 6.71 & 7.55 & 22\% & 6.65 & 6.65 & 0\% \\
128 & 5.52 & 5.16 & 76\% & \good{7.36} & 7.58 & 6\% \\
384 & \slow{4.73} & 4.27 & 86\% & 6.65 & 7.21 & 46\% \\
\bottomrule
\end{tabular}
\caption{B3 against \texttt{nc}. \#35, \texttt{jbpi}, core 4, uninstrumented, GOPs, average of three launches per
packer in A-B-B-A order. ``Slow reps'' is the share of repetitions in the slower of the two speeds.}
\label{tab:nc}
\end{table}

\paragraph{Step 4: the L1 prefetcher backs off (\#41, \#42).} The standard Linux cache counters read zero on the
X60, so \#42 found the X60's own event codes and added per-rep counts of L2 and prefetch activity. Comparing slow
reps with fast reps of the same run (instrumented build) gave a clear picture:
\begin{itemize}
  \item The L2 is accessed exactly as often in slow reps (ratio 0.95--1.01). The L2 is not overloaded.
  \item Demand misses in the L2 are 2--4 times higher in slow reps.
  \item Hits from the L1 data prefetcher fall to about half (ratio 0.40--0.56).
\end{itemize}
In other words, in a slow rep the L1 prefetcher does about half its usual work, and each line it no longer
brings in turns into a wait for memory. The same pattern appears with both packers, for B3 with K=256, and for
Reuse24's few slow reps. Running on core 2, which belongs to the other group of cores and has its own L2, gives
the same result, and no other core was busy during the timed runs. So the cause is not the shared L2 and not
another program. What makes the prefetcher back off is still unknown.

\paragraph{Decision on \texttt{nc} for B3.} After \#35 the hand-set value \texttt{nc} = 128 for B3 was first
removed, because the model cannot produce it. After \#41 and \#42 named the prefetcher as the cause, it was put
back as the baseline by the user's decision. B3 at \texttt{nc} = 128 still has two speeds with the vector packer.

\subsection{A first fix: clearing the pack lines (\#43)}
The X60 has no software prefetch instruction, but it does have \texttt{cbo.zero}, which sets a whole 64-byte
cache line to zero. The idea in \#43 was to clear each pack-buffer line just before packing into it, so that the
packer's writes would not first have to read the old line from the L2. The switch
(\texttt{--zero-pack-lines}) is off by default. It was measured over 168 launches: seven rounds on each board,
each round in the order off, on, on, off, off, on (Table~\ref{tab:cbo}).

\begin{table}[h]
\centering\small
\begin{tabular}{@{}llcccc@{}}
\toprule
& & \multicolumn{2}{c}{\textbf{by mean time}} & \multicolumn{2}{c}{\textbf{by median rep}} \\
\textbf{Case} & \textbf{Board} & off & on & off & on \\
\midrule
B3, \texttt{nc}=128 (baseline) & \texttt{jbpi}  & 5.71 & \good{6.04} & 5.42 & \good{6.20} \\
B3, \texttt{nc}=128 (baseline) & \texttt{jbpi2} & 5.56 & \good{6.01} & 5.01 & \good{5.98} \\
B3, \texttt{nc}=384 & \texttt{jbpi}  & 4.49 & 4.47 & 4.28 & 4.24 \\
B3, \texttt{nc}=384 & \texttt{jbpi2} & 4.64 & 4.47 & 4.31 & 4.33 \\
Reuse24 & \texttt{jbpi}  & 8.02 & \good{8.39} & 8.23 & \good{8.86} \\
Reuse24 & \texttt{jbpi2} & 7.96 & \good{8.32} & 8.14 & \good{8.64} \\
B2 & \texttt{jbpi}  & 12.40 & \slow{12.23} & 12.48 & 12.38 \\
B2 & \texttt{jbpi2} & 12.41 & \slow{12.27} & 12.48 & 12.37 \\
\bottomrule
\end{tabular}
\caption{The \texttt{cbo.zero} switch. \#43, core 4, uninstrumented, GOPs, average of 21 launches per side per
board.}
\label{tab:cbo}
\end{table}

The switch helps B3 at \texttt{nc} = 128 and Reuse24, costs B2 a small but steady amount (lower in all 14
rounds), and does nothing for B3 at \texttt{nc} = 384. However, the hardware counts show that it does
\emph{not} save any L2 reads: the packer's L2 writes rise by about one per cleared line, and its L2 reads do not
fall. So the idea behind the fix is not what makes it work. The gain is real but not yet explained.

\subsection{How reliable the measurements are (\#31, \#34, \#39)}
\label{sec:noise}
Several early findings turned out to be smaller than the noise in the measurements:
\begin{itemize}
  \item \textbf{Core 0 is slower.} It handles the board's device interrupts and runs the same program 5--10\%
  slower than core 4. All 32 launches of \#32 ran on core 0.
  \item \textbf{Launches differ.} The same program on the same core within the same hour varied by up to 1.7\%
  (\#33). Two cores in \#34 differed by up to 4\%, enough to reverse which tile won.
  \item \textbf{B3 is noisy inside a launch.} Its spread from rep to rep was 14--43\%, against about 1\% for the
  large shapes.
  \item \textbf{Some logs were not kept.} The core-0 swap runs of \#31 and the reference-side runs of \#33 left no
  logs, so their numbers cannot be checked again.
\end{itemize}
In response, \#39 set four house rules, now written in \texttt{docs/agents/execution.md}: never fix a benchmark
to core 0; treat one launch as no result, and compare small gaps inside one launch, repeated over at least three
launches in A-B-B-A order; quote only numbers from committed logs; and always name the build, timing method and
core. Issues \#35 onwards follow these rules.

%-------------------------------------------------------------------- 4
\section{Discussion}

\subsection{What explains the difference in speed}
Taken together, the issues remove the most obvious suspect and replace it with three smaller, more specific
causes. The code generator was the first suspect, but every direct test (register counts, spill counts, swap tests)
shows our kernels to be equal to the reference's. What does limit speed is:
\begin{enumerate}
  \item \textbf{Register budget.} Tiles that need more than 32 registers lose about 40\% of their speed on both
  designs. This is a property of the tile, not of the tool, and our shipped tiles stay within the budget.
  \item \textbf{Layer shape.} A narrow N (384) means more packing per operation. This explains why B4 trails the
  square shapes. The reference packs in the same way, so this is not a difference between the designs.
  \item \textbf{A memory slow state on small-M layers.} B3 has only 64 rows, so each packed block of B is used
  for little work, and memory behaviour dominates. In that situation the X60's L1 prefetcher can fall back to
  about half its work, and B3's speed halves with it.
\end{enumerate}
The third cause is the only one still open, and it is also the largest: B3 moves between roughly 4.5 and 7.5
GOPs depending on the packer, \texttt{nc} and time.

\subsection{Limitations}
\begin{itemize}
  \item These issues did not measure whether the \emph{reference} design's B3 also has two speeds. Since the
  cause lies in the processor's prefetcher and in the access pattern that both designs share, it probably does,
  but this has not been tested.
  \item The \#32 tile sweep ran entirely on core 0, and several numbers in \#31 and \#33 come from a single launch
  or have no saved log. The conclusions drawn from them hold because later, better-controlled runs agree, not
  because of those runs alone.
  \item The hardware-counter study of \#42 ran only on \texttt{jbpi}. The \#43 speed results cover both boards.
  \item The instrumented build changes timing slightly. Its counts explain behaviour, but its speed figures must
  not be compared with the reference.
  \item Why \texttt{cbo.zero} helps is not known. Turning it on by default without that knowledge risks keeping a
  gain that depends on timing by chance.
\end{itemize}

%-------------------------------------------------------------------- 5
\section{Conclusions}
\begin{enumerate}
  \item Our xDSL-generated microkernels are as good as the reference's Exo kernels: same registers, same spills,
  same speed within noise.
  \item Speed drops sharply only for tiles that need more than 32 vector registers. The shipped tiles (16x15 and
  32x7) are at or near the best; 32x6 and 32x7 are level for B4.
  \item B4 is slower than the square shapes because its narrow N forces about 2.3 times more packing per
  operation. Widening N to 1024 raises it from 8.10 to 9.02 GOPs.
  \item B3's instability is a memory slow state in which the L1 data prefetcher does about half its work. It is
  not the shared L2, not another program, and not the packer's own speed.
  \item Clearing the pack lines with \texttt{cbo.zero} raises B3 (\texttt{nc} = 128) from about 5.6 to about 6.0
  GOPs and Reuse24 from about 8.0 to 8.35 GOPs, lowers B2 from 12.40 to about 12.25 GOPs, and its reason is
  unknown.
  \item Differences of less than about 3\% between separate launches, and any figure measured on core 0, are not
  reliable. The \#39 house rules address this.
\end{enumerate}

%-------------------------------------------------------------------- 6
\section{Recommendations}
\begin{enumerate}
  \item \textbf{Keep the shipped tiles}, including 32x7 for B4, and keep \texttt{nc} = 128 as B3's baseline,
  always labelled as a hand-set value.
  \item \textbf{Run a control for \#43} with the same extra loop but without \texttt{cbo.zero}, before deciding
  whether to turn the switch on by default.
  \item \textbf{Measure the reference's B3 with the per-rep records}, to find out whether it also has two speeds.
  This decides whether B3's instability is a difference between the designs at all.
  \item \textbf{Continue \#41}: record the counters in short windows around the moment of the switch to the slow
  state, test padding the B row length away from 384, and repeat the counter study on \texttt{jbpi2}.
  \item \textbf{When the above is settled, consider \#44} (packing B one column at a time, so the prefetcher
  follows one stream instead of fifteen) and \#38 (the B4 packing costs).
  \item \textbf{Keep following the \#39 house rules} for every number that is quoted.
\end{enumerate}

%-------------------------------------------------------------------- references
\begin{thebibliography}{9}
\bibitem{b3134} \emph{Does packing or tile size limit performance? Brief and audit of issues \#31--\#34},
\texttt{.scratch/reports/issues-31-34-brief.pdf}, 27 September 2026.
\bibitem{b353639} \emph{Why B3 has two speeds, and what a smaller \texttt{nc} does. Brief of issues \#35, \#36 and
\#39}, \texttt{.scratch/reports/issues-35-36-39-brief.pdf}, 27 September 2026.
\bibitem{b42} \emph{Counting the X60's L2 and prefetcher per rep: what B3's slow state is. Brief of issue \#42},
\texttt{.scratch/reports/issue-42-brief.pdf}, 27 September 2026.
\bibitem{b43} \emph{Clearing the pack lines with \texttt{cbo.zero}: speed over 14 rounds on two boards. Brief of
issue \#43}, \texttt{.scratch/reports/issue-43-brief.pdf}, 28 September 2026.
\bibitem{issues} Issue files \texttt{.scratch/issues/0016}--\texttt{0028} and GitHub issues \#31--\#44,
\url{https://github.com/JieGH/RVV_code_gen_via_MLIR_xDSL/issues}.
\end{thebibliography}

%-------------------------------------------------------------------- appendix
\appendix
\newpage
\section{Issue-by-issue summary}
{\small
\begin{tabularx}{\linewidth}{@{}l >{\raggedright\arraybackslash}X >{\raggedright\arraybackslash}p{5.2cm}@{}}
\toprule
\textbf{Issue} & \textbf{Main result} & \textbf{Reliability} \\
\midrule
\#31 & The 32x7 tile is not B4's problem; our kernel and the reference's run level on core 4. & Conclusion holds; its
first evidence was indirect, and a claim about the k-loop was wrong (every case uses the same \texttt{kc}). \\
\#32 & Speed falls only above 32 registers; scalar packer wins on B3 only. & Holds; all runs on core 0; ``32x6 is
faster'' not proven. \\
\#33 & Same registers, identical spills, same speed as the reference at 16x16 and 24x10. & Holds; small
percentages were over-read; some logs not kept. \\
\#34 & B4's loss is packing per operation, caused by N=384; M a power of two adds a little. & Summarised, not
audited. \\
\#35 & Small \texttt{nc} steadies B3 only with the scalar packer; the model cannot produce it. & Three launches
per packer, house rules followed. \\
\#36 & B3's slowness is a memory slow state affecting packing and compute alike. & Instrumented build; one board.
\\
\#37 & (Plan) settle 32x6 against 32x7 inside one launch. & On hold. \\
\#38 & (Plan) cheaper A pack; pad row lengths off powers of two. & On hold. \\
\#39 & House rules for benchmarks; per-rep records. & Adopted. \\
\#41 & The L1 prefetcher does about half its work in slow reps; not the L2, not a neighbour. & Nine launches, one
board. \\
\#42 & X60 raw L2 and prefetch counters found, checked and recorded per rep. & 266 tests pass; reviewed. \\
\#43 & \texttt{cbo.zero} raises B3 (\texttt{nc}=128) and Reuse24, lowers B2 slightly; reason unknown. & 168
launches, two boards. \\
\#44 & (Plan) pack B one column at a time. & On hold. \\
\bottomrule
\end{tabularx}}

\end{document}
